\subsubsection{先行研究結果の定量的再現}
\label{sec:discussion_reproduction}

本節では、ペルー・ブラジル領域（ANC, HUA, EUS）を用いた解析結果と先行研究の結果を比較する。
第\ref{sec:results}章で示したように、ペルー・ブラジル領域における未発達型において、
南半球の冬至（J-solstice）に発生率が高く、夏至（D-solstice）に低いという明確な季節依存性を示した。
これは、\citeauthor{Ikesue2024}\citetype{Ikesue2024} (\citeyear{Ikesue2024}) で報告されている傾向と一致した。
また、潮汐効果との対応関係についても、
\citeauthor{Fujimoto2019} \citeyear{Fujimoto2019}で報告されているLunar age=6--9, 
18--21でDipで西向き電流が卓越するため、Dip EUELの値が下がり、発生率が高くなるという結果と一致した。
これらの結果は、本研究で開発した自動判定手法の妥当性を支持するものである。
